% Do not change document class, margins, fonts, etc.
\documentclass[a4paper,oneside,bibliography=totoc]{scrbook}

% some useful packages (add more as needed)
\usepackage[utf8]{inputenc}
\usepackage{graphicx}
\usepackage{latexsym}
\usepackage{amsmath}
\usepackage{amssymb}
\usepackage{tabularx}
\usepackage{booktabs}
\usepackage{algorithm} % you can modify the algorithm style to your liking
\counterwithin{algorithm}{chapter} % so that algorithms have chapter numbers as well
\usepackage{algorithmic}
\usepackage{csquotes}
\usepackage{comment}
\renewcommand{\algorithmiccomment}[1]{\hfill\textit{// #1}}
\usepackage[usenames,dvipsnames]{xcolor}
\usepackage[colorlinks,citecolor=Green]{hyperref} % you may change/remove the colors
\usepackage{lipsum} % you do not need this

\usepackage{natbib}
\bibliographystyle{chicagoa}
\setcitestyle{authoryear,round,semicolon,aysep={},yysep={,}} \let\cite\citep

% example enviroments (add more as needed)
\newtheorem{definition}{Definition} \newtheorem{proposition}{Proposition}

\begin{document}

\frontmatter \subject{Student Project} % change to appropriate type
\title{Flight Delay Prediction}
\author{Group 7 - Data Nuggets\\
  Cora Wackermann\\
  Kejsi Limaci\\
  Simon Schumacher\\
  Justus Mrosk\\
  Nick Birkholz\\
  Juan Cordero\\} \date{April 12, 2026}
\publishers{{\small Submitted to}\\
  Data and Web Science Group\\
  Prof. Christian Bizer\\
  University of Mannheim\\}
\maketitle
\begin{comment}
    \chapter{Abstract}

Your thesis must contain an abstract. A good reference for thesis writing
is~\citet{zobel2014}; we highly recommend that you study this or a similar book
during your studies. He writes the following about the abstract:

\blockcquote{zobel2004}{%
  An abstract is typically a single paragraph of about 50 to 200 words. The
  function of an abstract is to allow readers to judge whether or not the paper
  is of relevance to them. It should therefore be a concise summary of the
  paper's aims, scope, and conclusions. There is no space for unnecessary text;
  an abstract should be kept to as few words as possible while remaining clear
  and informative. Irrelevancies, such as minor details or a description of the
  structure of the paper, are inappropriate, as are acronyms, abbreviations, and
  mathematics. Sentences such as "We review relevant literature" should be
  omitted.

  The more specific an abstract is, the more interesting it is likely to be.
  Instead of writing "space requirements can be significantly reduced", write
  "space requirements can be reduced by 60\%". Instead of writing "we have a new
  inversion algorithm", write "we have a new inversion algorithm, based on
  move­to-front lists".

  Many scientists browse research papers outside their area of expertise. You
  should not assume that all likely readers will be specialists in the topic of
  their paper-abstracts should be self-contained and written for as broad a
  readership as possible. Only in rare circumstances should an abstract cite
  another paper (for example, when one paper consists entirely of analysis of
  results in another), in which case the reference should be given in full, not
  as a citation to the bibliography.}

% table of contents
\begingroup%
\hypersetup{hidelinks}% disable link color in TOC only
\tableofcontents%
\endgroup

% none of the things below is needed, but you may add them if you feel that they
% are helpful for your work \listofalgorithms \listoffigures \listoftables
% \listtheorems{definition} \listtheorems{proposition}
\end{comment} 


% okay, start new numbering ... here is where it really starts
\mainmatter

\chapter{Project Outline}
\label{ch:intro}

\section{What is the problem we are solving?}
In the aviation industry, operational efficiency is inextricably linked to the predictability of flight schedules. Flight delays represent more than just a passenger inconvenience; they impose significant financial burdens on airlines and create complex logistical bottlenecks for airport authorities. The goal of this data mining project is to develop a predictive framework to estimate flight delays with high precision.

The structure of this example thesis is \emph{non-binding}, but is highly
topic-dependent and should be discussed with your advisor. You probably do not
want to go below the level of subsections, though.
- Julains Checken therefeore maybe extende text 
\section{Data}
Our Project is mainly based on the "Carrier On-Time Performance" Data, from the American Buro of Transportation statistics, which can be downloaded from their website \link{https://www.transtats.bts.gov/Tables.asp?QO_VQ=EFD&QO_anzr=Nv4yv0r%FDb0-gvzr%FDcr4s14zn0pr%FDQn6n&QO_fu146_anzr=b0-gvzr}.
IT contains various features, about the flights, such as airline, airports, aircraft tail number and flight distance and especially information on their on time performance such as planned and actual departure and arrival times.
To keep the task feasible, we only use data from 2013 to 2019, to limit the amount of data and avoid year effected by covid19.
Additionally, we limit ourselves to 35 Mayor Airports in the US, and 7 mayor airlines, operating during that timeframe.
This makes the data more consistent and removes noise, while mostly removing less popular flights.
We additionally add historical weather measurement data by airport weather stations.
For that we use the meteostat python library, which provides hourly measurements for weather parameters like temperature, wind speed and precipitation.
This data originally stems from the National Oceanic and Atmospheric Administration (NOAA)\link{https://www.noaa.gov/}.
Beyond that we plan on experimenting, using additional pieces of information such as holidays, to build some features. 

\section{Approach to Solving the Problem}
- BUild Basline MOdel, whihc works by usign avergae arrival delays for arrival airplanse at arrival delay that day so far before the subject flight. 
 - we engeneer features, relavant to the issue.
    - by aggregeting data, and combine and reformulate the data
 - analyze relavant correltions
 - train differenet models on the Data, 
\subsection{Preprocessing steps}
 - Normalize/scale all numerical features and OneHotencode categorical features. (controll dimensionality and reduce nose)
- we throw out Diverserted anbd Cancelled flights, to have cleaner Data, more specific to our Task.
- we handle missing Values in approprate amnners depending on the feature.
- drop unsesary columns


\subsection{Algorithms to Be Utilized}
- Walk Forward Validation to avoid Data Leakage
- USe Last Year as Test Set.
- Needs to be like that to avoid Data Leakege of using Data later in time to train then the test data.
- hyperparameter Tuning to find best hyperpaprameter configs.
- Try out multiple models:
    - (Gradient Boosting Descsion Tree, Ensbles, Linear Regession, SVM, Bayes Regessors, and basic Neural network)
\section{Measuring success}
 - Mean Sqared Error
 - R^2 
 - 

\section{Goals}
- We want to have a significatly better predictions as our Basline Model
- we expectmodel to acchive this goal
- especially we expect Gradient Boosting to Perform well on the task.
- While Linear regession probably wont perform very well.
- We further expect to learn what factors ave high measureable influence on flight delays. 

\end{document}
